% Options for packages loaded elsewhere
% Options for packages loaded elsewhere
\PassOptionsToPackage{unicode}{hyperref}
\PassOptionsToPackage{hyphens}{url}
\PassOptionsToPackage{dvipsnames,svgnames,x11names}{xcolor}
%
\documentclass[
  letterpaper,
  DIV=11,
  numbers=noendperiod]{scrartcl}
\usepackage{xcolor}
\usepackage{amsmath,amssymb}
\setcounter{secnumdepth}{-\maxdimen} % remove section numbering
\usepackage{iftex}
\ifPDFTeX
  \usepackage[T1]{fontenc}
  \usepackage[utf8]{inputenc}
  \usepackage{textcomp} % provide euro and other symbols
\else % if luatex or xetex
  \usepackage{unicode-math} % this also loads fontspec
  \defaultfontfeatures{Scale=MatchLowercase}
  \defaultfontfeatures[\rmfamily]{Ligatures=TeX,Scale=1}
\fi
\usepackage{lmodern}
\ifPDFTeX\else
  % xetex/luatex font selection
\fi
% Use upquote if available, for straight quotes in verbatim environments
\IfFileExists{upquote.sty}{\usepackage{upquote}}{}
\IfFileExists{microtype.sty}{% use microtype if available
  \usepackage[]{microtype}
  \UseMicrotypeSet[protrusion]{basicmath} % disable protrusion for tt fonts
}{}
\makeatletter
\@ifundefined{KOMAClassName}{% if non-KOMA class
  \IfFileExists{parskip.sty}{%
    \usepackage{parskip}
  }{% else
    \setlength{\parindent}{0pt}
    \setlength{\parskip}{6pt plus 2pt minus 1pt}}
}{% if KOMA class
  \KOMAoptions{parskip=half}}
\makeatother
% Make \paragraph and \subparagraph free-standing
\makeatletter
\ifx\paragraph\undefined\else
  \let\oldparagraph\paragraph
  \renewcommand{\paragraph}{
    \@ifstar
      \xxxParagraphStar
      \xxxParagraphNoStar
  }
  \newcommand{\xxxParagraphStar}[1]{\oldparagraph*{#1}\mbox{}}
  \newcommand{\xxxParagraphNoStar}[1]{\oldparagraph{#1}\mbox{}}
\fi
\ifx\subparagraph\undefined\else
  \let\oldsubparagraph\subparagraph
  \renewcommand{\subparagraph}{
    \@ifstar
      \xxxSubParagraphStar
      \xxxSubParagraphNoStar
  }
  \newcommand{\xxxSubParagraphStar}[1]{\oldsubparagraph*{#1}\mbox{}}
  \newcommand{\xxxSubParagraphNoStar}[1]{\oldsubparagraph{#1}\mbox{}}
\fi
\makeatother

\usepackage{color}
\usepackage{fancyvrb}
\newcommand{\VerbBar}{|}
\newcommand{\VERB}{\Verb[commandchars=\\\{\}]}
\DefineVerbatimEnvironment{Highlighting}{Verbatim}{commandchars=\\\{\}}
% Add ',fontsize=\small' for more characters per line
\usepackage{framed}
\definecolor{shadecolor}{RGB}{241,243,245}
\newenvironment{Shaded}{\begin{snugshade}}{\end{snugshade}}
\newcommand{\AlertTok}[1]{\textcolor[rgb]{0.68,0.00,0.00}{#1}}
\newcommand{\AnnotationTok}[1]{\textcolor[rgb]{0.37,0.37,0.37}{#1}}
\newcommand{\AttributeTok}[1]{\textcolor[rgb]{0.40,0.45,0.13}{#1}}
\newcommand{\BaseNTok}[1]{\textcolor[rgb]{0.68,0.00,0.00}{#1}}
\newcommand{\BuiltInTok}[1]{\textcolor[rgb]{0.00,0.23,0.31}{#1}}
\newcommand{\CharTok}[1]{\textcolor[rgb]{0.13,0.47,0.30}{#1}}
\newcommand{\CommentTok}[1]{\textcolor[rgb]{0.37,0.37,0.37}{#1}}
\newcommand{\CommentVarTok}[1]{\textcolor[rgb]{0.37,0.37,0.37}{\textit{#1}}}
\newcommand{\ConstantTok}[1]{\textcolor[rgb]{0.56,0.35,0.01}{#1}}
\newcommand{\ControlFlowTok}[1]{\textcolor[rgb]{0.00,0.23,0.31}{\textbf{#1}}}
\newcommand{\DataTypeTok}[1]{\textcolor[rgb]{0.68,0.00,0.00}{#1}}
\newcommand{\DecValTok}[1]{\textcolor[rgb]{0.68,0.00,0.00}{#1}}
\newcommand{\DocumentationTok}[1]{\textcolor[rgb]{0.37,0.37,0.37}{\textit{#1}}}
\newcommand{\ErrorTok}[1]{\textcolor[rgb]{0.68,0.00,0.00}{#1}}
\newcommand{\ExtensionTok}[1]{\textcolor[rgb]{0.00,0.23,0.31}{#1}}
\newcommand{\FloatTok}[1]{\textcolor[rgb]{0.68,0.00,0.00}{#1}}
\newcommand{\FunctionTok}[1]{\textcolor[rgb]{0.28,0.35,0.67}{#1}}
\newcommand{\ImportTok}[1]{\textcolor[rgb]{0.00,0.46,0.62}{#1}}
\newcommand{\InformationTok}[1]{\textcolor[rgb]{0.37,0.37,0.37}{#1}}
\newcommand{\KeywordTok}[1]{\textcolor[rgb]{0.00,0.23,0.31}{\textbf{#1}}}
\newcommand{\NormalTok}[1]{\textcolor[rgb]{0.00,0.23,0.31}{#1}}
\newcommand{\OperatorTok}[1]{\textcolor[rgb]{0.37,0.37,0.37}{#1}}
\newcommand{\OtherTok}[1]{\textcolor[rgb]{0.00,0.23,0.31}{#1}}
\newcommand{\PreprocessorTok}[1]{\textcolor[rgb]{0.68,0.00,0.00}{#1}}
\newcommand{\RegionMarkerTok}[1]{\textcolor[rgb]{0.00,0.23,0.31}{#1}}
\newcommand{\SpecialCharTok}[1]{\textcolor[rgb]{0.37,0.37,0.37}{#1}}
\newcommand{\SpecialStringTok}[1]{\textcolor[rgb]{0.13,0.47,0.30}{#1}}
\newcommand{\StringTok}[1]{\textcolor[rgb]{0.13,0.47,0.30}{#1}}
\newcommand{\VariableTok}[1]{\textcolor[rgb]{0.07,0.07,0.07}{#1}}
\newcommand{\VerbatimStringTok}[1]{\textcolor[rgb]{0.13,0.47,0.30}{#1}}
\newcommand{\WarningTok}[1]{\textcolor[rgb]{0.37,0.37,0.37}{\textit{#1}}}

\usepackage{longtable,booktabs,array}
\usepackage{calc} % for calculating minipage widths
% Correct order of tables after \paragraph or \subparagraph
\usepackage{etoolbox}
\makeatletter
\patchcmd\longtable{\par}{\if@noskipsec\mbox{}\fi\par}{}{}
\makeatother
% Allow footnotes in longtable head/foot
\IfFileExists{footnotehyper.sty}{\usepackage{footnotehyper}}{\usepackage{footnote}}
\makesavenoteenv{longtable}
\usepackage{graphicx}
\makeatletter
\newsavebox\pandoc@box
\newcommand*\pandocbounded[1]{% scales image to fit in text height/width
  \sbox\pandoc@box{#1}%
  \Gscale@div\@tempa{\textheight}{\dimexpr\ht\pandoc@box+\dp\pandoc@box\relax}%
  \Gscale@div\@tempb{\linewidth}{\wd\pandoc@box}%
  \ifdim\@tempb\p@<\@tempa\p@\let\@tempa\@tempb\fi% select the smaller of both
  \ifdim\@tempa\p@<\p@\scalebox{\@tempa}{\usebox\pandoc@box}%
  \else\usebox{\pandoc@box}%
  \fi%
}
% Set default figure placement to htbp
\def\fps@figure{htbp}
\makeatother





\setlength{\emergencystretch}{3em} % prevent overfull lines

\providecommand{\tightlist}{%
  \setlength{\itemsep}{0pt}\setlength{\parskip}{0pt}}



 


\KOMAoption{captions}{tableheading}
\makeatletter
\@ifpackageloaded{caption}{}{\usepackage{caption}}
\AtBeginDocument{%
\ifdefined\contentsname
  \renewcommand*\contentsname{Table of contents}
\else
  \newcommand\contentsname{Table of contents}
\fi
\ifdefined\listfigurename
  \renewcommand*\listfigurename{List of Figures}
\else
  \newcommand\listfigurename{List of Figures}
\fi
\ifdefined\listtablename
  \renewcommand*\listtablename{List of Tables}
\else
  \newcommand\listtablename{List of Tables}
\fi
\ifdefined\figurename
  \renewcommand*\figurename{Figure}
\else
  \newcommand\figurename{Figure}
\fi
\ifdefined\tablename
  \renewcommand*\tablename{Table}
\else
  \newcommand\tablename{Table}
\fi
}
\@ifpackageloaded{float}{}{\usepackage{float}}
\floatstyle{ruled}
\@ifundefined{c@chapter}{\newfloat{codelisting}{h}{lop}}{\newfloat{codelisting}{h}{lop}[chapter]}
\floatname{codelisting}{Listing}
\newcommand*\listoflistings{\listof{codelisting}{List of Listings}}
\makeatother
\makeatletter
\makeatother
\makeatletter
\@ifpackageloaded{caption}{}{\usepackage{caption}}
\@ifpackageloaded{subcaption}{}{\usepackage{subcaption}}
\makeatother
\usepackage{bookmark}
\IfFileExists{xurl.sty}{\usepackage{xurl}}{} % add URL line breaks if available
\urlstyle{same}
\hypersetup{
  pdftitle={SImulation plan},
  colorlinks=true,
  linkcolor={blue},
  filecolor={Maroon},
  citecolor={Blue},
  urlcolor={Blue},
  pdfcreator={LaTeX via pandoc}}


\title{SImulation plan}
\author{}
\date{}
\begin{document}
\maketitle


\section{Simulation Plan for Bayesian Borrowing with Time-Dependent
Historical
Controls}\label{simulation-plan-for-bayesian-borrowing-with-time-dependent-historical-controls}

\subsection{1. Purpose of the Simulation
Study}\label{purpose-of-the-simulation-study}

This document sets out a simulation plan for evaluating Bayesian
borrowing methods for a randomized clinical trial with a continuous
endpoint when patient-level data are available from multiple historical
control studies. The focus is on settings where the historical control
mean may drift over time. The simulation will compare methods that
ignore drift, partially account for drift, or explicitly model temporal
dependence.

The main aims are:

\begin{enumerate}
\def\labelenumi{\arabic{enumi}.}
\tightlist
\item
  To assess bias, precision, interval estimation, and decision
  performance for the treatment effect when borrowing historical
  information.
\item
  To evaluate robustness of competing borrowing methods under no drift,
  mild drift, and substantial drift.
\item
  To study the effect of covariate imbalance and outcome-model
  misspecification on borrowing performance.
\item
  To estimate marginal treatment and control means using g-computation
  and then obtain the marginal treatment effect.
\end{enumerate}

\begin{center}\rule{0.5\linewidth}{0.5pt}\end{center}

\subsection{2. Overall Trial Setting}\label{overall-trial-setting}

We consider:

\begin{itemize}
\tightlist
\item
  \textbf{H = 4 historical studies}, each containing only control-arm
  patients.
\item
  \textbf{One current randomized trial} with both treatment and control
  arms.
\item
  \textbf{Continuous outcome}.
\item
  \textbf{Patient-level baseline covariates} available for all studies.
\item
  The estimand is defined for the \textbf{current trial population}.
\end{itemize}

Let:

\begin{itemize}
\tightlist
\item
  (h = 1,\dots,H) index the historical studies.
\item
  (s = H+1) denote the current study.
\item
  (i) index individuals within a study.
\item
  (A\_\{si\} \in \{0,1\}) be treatment assignment in the current trial,
  where 1 = treatment and 0 = control.
\item
  (X\_\{si\}) be the baseline covariate vector.
\item
  (Y\_\{si\}) be the continuous outcome.
\end{itemize}

The historical studies contain only ((Y\_\{hi\}, X\_\{hi\})) for
controls.

\begin{center}\rule{0.5\linewidth}{0.5pt}\end{center}

\subsection{3. Target Estimands}\label{target-estimands}

The estimands are marginal means under the covariate distribution of the
\textbf{current trial population}:

{[} \mu\_T = E\{Y(1) \mid S = H+1\}, \qquad \mu\_C = E\{Y(0) \mid S =
H+1\}, {]}

and the marginal treatment effect:

{[} \Delta = \mu\_T - \mu\_C. {]}

These will be estimated using posterior g-computation:

{[} \mu\emph{T = \frac{1}{n_{\text{cur}}}
\sum}\{i=1\}\^{}\{n\_\{\text{cur}\}\} E(Y \mid A=1, X=X\_i, \Theta), {]}

{[} \mu\emph{C = \frac{1}{n_{\text{cur}}}
\sum}\{i=1\}\^{}\{n\_\{\text{cur}\}\} E(Y \mid A=0, X=X\_i, \Theta), {]}

where (\Theta) denotes the model parameters and the averaging is over
the observed covariate distribution in the current trial.

\begin{center}\rule{0.5\linewidth}{0.5pt}\end{center}

\subsection{4. Data-Generating
Mechanism}\label{data-generating-mechanism}

\subsubsection{4.1 Baseline Covariates}\label{baseline-covariates}

To reflect a realistic mixture of covariates, let (p=4):

\begin{itemize}
\tightlist
\item
  (X\_1 \sim N(0,1))
\item
  (X\_2 \sim \text{Bernoulli}(0.5))
\item
  (X\_3 \sim N(0,1))
\item
  (X\_4 \sim \text{Bernoulli}(0.4))
\end{itemize}

independently unless otherwise stated.

A sensitivity analysis may replace the independent structure with
correlated continuous covariates:

{[} (X\_1, X\_3)\^{}\top \sim N\_2\left(

\begin{pmatrix}0\\0\end{pmatrix}

,

\begin{pmatrix}1 & \rho \\\rho & 1\end{pmatrix}

\right), {]}

with (\rho = 0.3).

\subsubsection{4.2 Outcome Model Used for Data
Generation}\label{outcome-model-used-for-data-generation}

For individual (i) in study (j), generate outcomes from

{[} Y\_\{ji\} = \alpha\emph{j + \tau A}\{ji\} +
X\_\{ji\}\^{}\top \beta + \varepsilon\emph{\{ji\},
\qquad \varepsilon}\{ji\} \sim N(0, \sigma\^{}2). {]}

Here:

\begin{itemize}
\tightlist
\item
  (\alpha\_j) is the study-specific intercept capturing the study-level
  control mean after covariate adjustment.
\item
  (\tau) is the treatment effect in the current trial only.
\item
  (\beta) are common prognostic effects.
\item
  (\sigma\^{}2) is the residual variance.
\end{itemize}

For historical studies, (A\_\{hi\}=0) for all individuals. For the
current trial, (A\_\{si\}\sim \text{Bernoulli}(0.5)).

\subsubsection{4.3 Parameterization of the True Outcome
Model}\label{parameterization-of-the-true-outcome-model}

A practical baseline parameterization is:

{[} \beta = (0.6, -0.5, 0.4, -0.3)\^{}\top, \qquad \sigma = 1. {]}

The true treatment effect may be set to:

\begin{itemize}
\tightlist
\item
  \textbf{Null scenario}: (\tau = 0)
\item
  \textbf{Beneficial treatment scenario}: (\tau = -0.3)
\item
  \textbf{Moderate beneficial effect}: (\tau = -0.5)
\end{itemize}

assuming smaller values are better. If larger values are better, the
signs can be reversed.

\subsubsection{4.4 Time-Dependent Control Mean
Structure}\label{time-dependent-control-mean-structure}

The historical studies are assumed ordered in time. Let
(\alpha\_1,\dots,\alpha\emph{H,\alpha}\{H+1\}) denote the adjusted
control intercepts for the historical and current studies.

\paragraph{Scenario A: No drift}\label{scenario-a-no-drift}

{[} \alpha\_j = 0 \quad \text{for all } j=1,\dots,H+1. {]}

\paragraph{Scenario B: Linear drift}\label{scenario-b-linear-drift}

{[} \alpha\_j = \delta\_0 + \delta\_1(j-1), {]}

with, for example, (\delta\_0=0) and (\delta\_1=0.15).

\paragraph{Scenario C: Step change near the current
trial}\label{scenario-c-step-change-near-the-current-trial}

{[} \alpha\_1 = 0,\quad \alpha\_2 = 0,\quad \alpha\_3 =
0.2,\quad \alpha\_4 = 0.4,\quad \alpha\_5 = 0.7. {]}

\paragraph{Scenario D: AR(1) drift}\label{scenario-d-ar1-drift}

{[} \alpha\emph{1 \sim N(0, \omega\textsuperscript{2/(1-\phi}2)), \qquad
\alpha\emph{j \mid \alpha}\{j-1\} \sim N(\phi \alpha}\{j-1\},
\omega\^{}2), ; j=2,\dots,H+1. {]}

Suggested values:

\begin{itemize}
\tightlist
\item
  Mild dependence: (\phi = 0.5), (\omega = 0.15)
\item
  Strong dependence: (\phi = 0.8), (\omega = 0.10)
\end{itemize}

This AR(1) structure will be especially useful when evaluating
time-dependent priors.

\begin{center}\rule{0.5\linewidth}{0.5pt}\end{center}

\subsection{5. Sample Size Settings}\label{sample-size-settings}

A default sample-size design is:

\begin{itemize}
\tightlist
\item
  Historical studies: (n\_h \in \{80,80,80,80\}) or (\{50,100,150,200\})
\item
  Current control: (n\_\{C,\text{cur}\} = 100)
\item
  Current treatment: (n\_\{T,\text{cur}\} = 100)
\end{itemize}

Sensitivity analyses may consider a smaller current trial, for example:

\begin{itemize}
\tightlist
\item
  Current control = 50
\item
  Current treatment = 50
\end{itemize}

This is important because borrowing is usually more influential when the
current trial is small.

\begin{center}\rule{0.5\linewidth}{0.5pt}\end{center}

\subsection{6. Competing Analysis
Methods}\label{competing-analysis-methods}

The following methods will be compared.

\subsubsection{6.1 No Borrowing}\label{no-borrowing}

Fit the model using only current-trial data:

{[} Y\_i = \alpha\_\{\text{cur}\} + \tau A\_i + X\_i\^{}\top \beta +
\varepsilon\_i. {]}

This provides the benchmark with no external information.

\subsubsection{6.2 Full Borrowing (Pooled
Controls)}\label{full-borrowing-pooled-controls}

Pool all historical controls with the current control under a common
control intercept:

{[} Y\_i = \alpha + \tau A\_i + X\_i\^{}\top \beta + \varepsilon\_i. {]}

This method is efficient under no drift but can perform poorly when
drift is present.

\subsubsection{6.3 Hierarchical Exchangeability
Model}\label{hierarchical-exchangeability-model}

Allow study-specific intercepts but assume exchangeability:

{[} \alpha\emph{j \sim N(\mu}\alpha, \tau\_\alpha\^{}2),
\qquad j=1,\dots,H+1. {]}

This permits partial borrowing across studies without explicit time
ordering.

\subsubsection{6.4 AR(1) Time-Dependent
Prior}\label{ar1-time-dependent-prior}

Assume the study-specific intercepts evolve over time:

{[} \alpha\emph{1 \sim N(m\_0, s\_0\^{}2), \qquad \alpha\emph{j
\mid \alpha}\{j-1\} \sim N(\phi \alpha}\{j-1\}, \omega\^{}2), ;
j=2,\dots,H+1. {]}

This is the main time-dependent borrowing model.

\subsubsection{6.5 Power Prior}\label{power-prior}

Use the historical likelihood contribution with downweighting parameters
(a\_h \in [0,1]):

{[} \pi(\Theta \mid D\_1,\dots,D\_H,D\_\{\text{cur}\}) \propto
\left\{\prod\emph{\{h=1\}\^{}\{H\}
L(\Theta \mid D\_h)\^{}\{a\_h\}\right\}
L(\Theta \mid D}\{\text{cur}\})\pi\_0(\Theta). {]}

Two implementations can be considered:

\begin{enumerate}
\def\labelenumi{\arabic{enumi}.}
\tightlist
\item
  \textbf{Fixed power prior}, e.g.~(a\_h = 0.5) or decreasing with age
  of the study.
\item
  \textbf{Normalized random power prior}, e.g.~(a\_h
  \sim \text{Beta}(c\_1,c\_2)).
\end{enumerate}

A simple time-sensitive choice is:

{[} a\_h = \exp\{-\lambda(H+1-h)\}, {]}

with (\lambda \textgreater{} 0), so older studies receive less weight.

\begin{center}\rule{0.5\linewidth}{0.5pt}\end{center}

\subsection{7. Bayesian Analysis Model Used for
G-Computation}\label{bayesian-analysis-model-used-for-g-computation}

A convenient unified analysis model is:

{[} Y\_\{ji\} \sim N(\alpha\emph{j + \tau , I(j=H+1)A}\{ji\} +
X\_\{ji\}\^{}\top \beta, \sigma\^{}2). {]}

The difference between methods lies mainly in the prior structure
assigned to (\alpha\_j).

\subsubsection{7.1 Prior Specification}\label{prior-specification}

A weakly informative prior specification is:

{[} \tau \sim N(0, 2\^{}2), \qquad \beta\_k \sim N(0, 2\^{}2),
\qquad \sigma \sim \text{half-Normal}(0, 1.5). {]}

For the hierarchical model:

{[} \mu\emph{\alpha \sim N(0, 2\^{}2), \qquad
\tau}\alpha \sim \text{half-Normal}(0,1). {]}

For the AR(1) model:

{[} m\_0 \sim N(0,2\^{}2), \qquad s\_0 \sim \text{half-Normal}(0,1),
\qquad \phi \sim \text{Uniform}(-1,1), \qquad
\omega \sim \text{half-Normal}(0,1). {]}

A more stable alternative for (\phi) in Stan is to estimate an
unconstrained parameter and transform it to ((-1,1)).

\begin{center}\rule{0.5\linewidth}{0.5pt}\end{center}

\subsection{8. Simulation Scenarios}\label{simulation-scenarios}

The simulation factors may include:

\begin{enumerate}
\def\labelenumi{\arabic{enumi}.}
\tightlist
\item
  \textbf{Drift pattern}

  \begin{itemize}
  \tightlist
  \item
    None
  \item
    Linear drift
  \item
    Step change
  \item
    AR(1) drift
  \end{itemize}
\item
  \textbf{Treatment effect}

  \begin{itemize}
  \tightlist
  \item
    (\tau = 0)
  \item
    (\tau = -0.3)
  \item
    (\tau = -0.5)
  \end{itemize}
\item
  \textbf{Current sample size}

  \begin{itemize}
  \tightlist
  \item
    100 + 100
  \item
    50 + 50
  \end{itemize}
\item
  \textbf{Historical sample-size balance}

  \begin{itemize}
  \tightlist
  \item
    Equal sizes
  \item
    Unequal sizes
  \end{itemize}
\item
  \textbf{Covariate distribution shift between historical and current
  studies}

  \begin{itemize}
  \tightlist
  \item
    None
  \item
    Mild shift, e.g.~(P(X\_2=1)) changes from 0.5 to 0.65 in the current
    trial
  \end{itemize}
\item
  \textbf{Outcome-model misspecification}

  \begin{itemize}
  \tightlist
  \item
    Correct linear model
  \item
    Omitted interaction, e.g.~(0.3X\_1X\_2) added in generation but not
    in fitting
  \end{itemize}
\end{enumerate}

\begin{center}\rule{0.5\linewidth}{0.5pt}\end{center}

\subsection{9. Performance Measures}\label{performance-measures}

For each method and scenario, evaluate:

\begin{enumerate}
\def\labelenumi{\arabic{enumi}.}
\tightlist
\item
  \textbf{Bias} of (\hat\mu\_T), (\hat\mu\_C), and (\hat\Delta)
\item
  \textbf{Empirical standard deviation}
\item
  \textbf{Mean posterior standard deviation}
\item
  \textbf{Root mean squared error (RMSE)}
\item
  \textbf{95\% credible interval coverage}
\item
  \textbf{Average credible interval width}
\item
  \textbf{Posterior probability of benefit}, e.g. {[}
  P(\Delta \textless{} 0 \mid \text{data}) {]} when smaller outcomes are
  better
\item
  \textbf{Type I error / false positive rate} under (\tau=0)
\item
  \textbf{Power} under (\tau \textless{} 0)
\end{enumerate}

A decision rule may be:

{[} \text{Declare efficacy if } P(\Delta \textless{} 0 \mid \text{data})
\textgreater{} 0.975. {]}

\begin{center}\rule{0.5\linewidth}{0.5pt}\end{center}

\subsection{10. Monte Carlo Design}\label{monte-carlo-design}

For each scenario:

\begin{enumerate}
\def\labelenumi{\arabic{enumi}.}
\tightlist
\item
  Generate covariates and outcomes for the 4 historical studies and the
  current trial.
\item
  Fit each Bayesian analysis model.
\item
  Draw posterior samples for model parameters.
\item
  Compute posterior g-computation estimates of (\mu\_T), (\mu\_C), and
  (\Delta).
\item
  Store point estimates, interval estimates, and decision indicators.
\item
  Repeat for (R) simulated datasets.
\end{enumerate}

A suitable starting choice is:

\begin{itemize}
\tightlist
\item
  (R = 1000) Monte Carlo replicates for final results
\item
  2 to 4 chains
\item
  1000 warmup + 1000 sampling iterations per chain
\end{itemize}

For pilot runs, use smaller (R) and fewer iterations.

\begin{center}\rule{0.5\linewidth}{0.5pt}\end{center}

\subsection{11. Suggested Stan
Parameterization}\label{suggested-stan-parameterization}

To improve computation, it is useful to center and scale continuous
covariates before fitting the models. For hierarchical and AR(1) models,
a non-centered parameterization is preferred.

\subsubsection{11.1 Non-centered hierarchical
intercepts}\label{non-centered-hierarchical-intercepts}

Instead of sampling

{[} \alpha\emph{j \sim N(\mu}\alpha, \tau\_\alpha\^{}2), {]}

sample

{[} z\_j \sim N(0,1), \qquad \alpha\emph{j = \mu}\alpha +
\tau\_\alpha z\_j. {]}

\subsubsection{11.2 Transformed AR(1)
coefficient}\label{transformed-ar1-coefficient}

Sample (\eta\_\phi \in \mathbb{R}) and define

{[} \phi = 2,\mathrm{inv\_logit}(\eta\_\phi)-1. {]}

In Stan, this can be written using \texttt{inv\_logit}.

\begin{center}\rule{0.5\linewidth}{0.5pt}\end{center}

\subsection{12. Stan Code: No Borrowing
Model}\label{stan-code-no-borrowing-model}

\begin{Shaded}
\begin{Highlighting}[]
\KeywordTok{data}\NormalTok{ \{}
  \DataTypeTok{int}\NormalTok{\textless{}}\KeywordTok{lower}\NormalTok{=}\DecValTok{1}\NormalTok{\textgreater{} N;}
  \DataTypeTok{int}\NormalTok{\textless{}}\KeywordTok{lower}\NormalTok{=}\DecValTok{1}\NormalTok{\textgreater{} P;}
  \DataTypeTok{matrix}\NormalTok{[N, P] X;}
  \DataTypeTok{vector}\NormalTok{[N] y;}
  \DataTypeTok{array}\NormalTok{[N] }\DataTypeTok{int}\NormalTok{\textless{}}\KeywordTok{lower}\NormalTok{=}\DecValTok{0}\NormalTok{, }\KeywordTok{upper}\NormalTok{=}\DecValTok{1}\NormalTok{\textgreater{} A;}
\NormalTok{\}}
\KeywordTok{parameters}\NormalTok{ \{}
  \DataTypeTok{real}\NormalTok{ alpha;}
  \DataTypeTok{real}\NormalTok{ tau;}
  \DataTypeTok{vector}\NormalTok{[P] beta;}
  \DataTypeTok{real}\NormalTok{\textless{}}\KeywordTok{lower}\NormalTok{=}\DecValTok{0}\NormalTok{\textgreater{} sigma;}
\NormalTok{\}}
\KeywordTok{model}\NormalTok{ \{}
\NormalTok{  alpha \textasciitilde{} normal(}\DecValTok{0}\NormalTok{, }\DecValTok{2}\NormalTok{);}
\NormalTok{  tau   \textasciitilde{} normal(}\DecValTok{0}\NormalTok{, }\DecValTok{2}\NormalTok{);}
\NormalTok{  beta  \textasciitilde{} normal(}\DecValTok{0}\NormalTok{, }\DecValTok{2}\NormalTok{);}
\NormalTok{  sigma \textasciitilde{} normal(}\DecValTok{0}\NormalTok{, }\FloatTok{1.5}\NormalTok{);}

\NormalTok{  y \textasciitilde{} normal(alpha + tau * to\_vector(A) + X * beta, sigma);}
\NormalTok{\}}
\KeywordTok{generated quantities}\NormalTok{ \{}
  \DataTypeTok{vector}\NormalTok{[N] mu1;}
  \DataTypeTok{vector}\NormalTok{[N] mu0;}
  \DataTypeTok{real}\NormalTok{ muT;}
  \DataTypeTok{real}\NormalTok{ muC;}
  \DataTypeTok{real}\NormalTok{ Delta;}

\NormalTok{  mu1 = alpha + tau * rep\_vector(}\DecValTok{1}\NormalTok{, N) + X * beta;}
\NormalTok{  mu0 = alpha + X * beta;}
\NormalTok{  muT = mean(mu1);}
\NormalTok{  muC = mean(mu0);}
\NormalTok{  Delta = muT {-} muC;}
\NormalTok{\}}
\end{Highlighting}
\end{Shaded}

\begin{center}\rule{0.5\linewidth}{0.5pt}\end{center}

\subsection{13. Stan Code: Full Borrowing (Pooled
Controls)}\label{stan-code-full-borrowing-pooled-controls}

\begin{Shaded}
\begin{Highlighting}[]
\KeywordTok{data}\NormalTok{ \{}
  \DataTypeTok{int}\NormalTok{\textless{}}\KeywordTok{lower}\NormalTok{=}\DecValTok{1}\NormalTok{\textgreater{} N;}
  \DataTypeTok{int}\NormalTok{\textless{}}\KeywordTok{lower}\NormalTok{=}\DecValTok{1}\NormalTok{\textgreater{} P;}
  \DataTypeTok{int}\NormalTok{\textless{}}\KeywordTok{lower}\NormalTok{=}\DecValTok{1}\NormalTok{\textgreater{} S;}
  \DataTypeTok{matrix}\NormalTok{[N, P] X;}
  \DataTypeTok{vector}\NormalTok{[N] y;}
  \DataTypeTok{array}\NormalTok{[N] }\DataTypeTok{int}\NormalTok{\textless{}}\KeywordTok{lower}\NormalTok{=}\DecValTok{1}\NormalTok{, }\KeywordTok{upper}\NormalTok{=S\textgreater{} study;}
  \DataTypeTok{array}\NormalTok{[N] }\DataTypeTok{int}\NormalTok{\textless{}}\KeywordTok{lower}\NormalTok{=}\DecValTok{0}\NormalTok{, }\KeywordTok{upper}\NormalTok{=}\DecValTok{1}\NormalTok{\textgreater{} current;}
  \DataTypeTok{array}\NormalTok{[N] }\DataTypeTok{int}\NormalTok{\textless{}}\KeywordTok{lower}\NormalTok{=}\DecValTok{0}\NormalTok{, }\KeywordTok{upper}\NormalTok{=}\DecValTok{1}\NormalTok{\textgreater{} A;}
\NormalTok{\}}
\KeywordTok{parameters}\NormalTok{ \{}
  \DataTypeTok{real}\NormalTok{ alpha;}
  \DataTypeTok{real}\NormalTok{ tau;}
  \DataTypeTok{vector}\NormalTok{[P] beta;}
  \DataTypeTok{real}\NormalTok{\textless{}}\KeywordTok{lower}\NormalTok{=}\DecValTok{0}\NormalTok{\textgreater{} sigma;}
\NormalTok{\}}
\KeywordTok{model}\NormalTok{ \{}
\NormalTok{  alpha \textasciitilde{} normal(}\DecValTok{0}\NormalTok{, }\DecValTok{2}\NormalTok{);}
\NormalTok{  tau   \textasciitilde{} normal(}\DecValTok{0}\NormalTok{, }\DecValTok{2}\NormalTok{);}
\NormalTok{  beta  \textasciitilde{} normal(}\DecValTok{0}\NormalTok{, }\DecValTok{2}\NormalTok{);}
\NormalTok{  sigma \textasciitilde{} normal(}\DecValTok{0}\NormalTok{, }\FloatTok{1.5}\NormalTok{);}

  \ControlFlowTok{for}\NormalTok{ (n }\ControlFlowTok{in} \DecValTok{1}\NormalTok{:N) \{}
\NormalTok{    y[n] \textasciitilde{} normal(alpha + tau * current[n] * A[n] + X[n] * beta, sigma);}
\NormalTok{  \}}
\NormalTok{\}}
\end{Highlighting}
\end{Shaded}

This model pools all controls through a common intercept \texttt{alpha}.

\begin{center}\rule{0.5\linewidth}{0.5pt}\end{center}

\subsection{14. Stan Code: Hierarchical Exchangeability
Model}\label{stan-code-hierarchical-exchangeability-model}

\begin{Shaded}
\begin{Highlighting}[]
\KeywordTok{data}\NormalTok{ \{}
  \DataTypeTok{int}\NormalTok{\textless{}}\KeywordTok{lower}\NormalTok{=}\DecValTok{1}\NormalTok{\textgreater{} N;}
  \DataTypeTok{int}\NormalTok{\textless{}}\KeywordTok{lower}\NormalTok{=}\DecValTok{1}\NormalTok{\textgreater{} P;}
  \DataTypeTok{int}\NormalTok{\textless{}}\KeywordTok{lower}\NormalTok{=}\DecValTok{1}\NormalTok{\textgreater{} S;                 }\CommentTok{// 4 historical + current = 5}
  \DataTypeTok{int}\NormalTok{\textless{}}\KeywordTok{lower}\NormalTok{=}\DecValTok{1}\NormalTok{\textgreater{} Ncur;              }\CommentTok{// current trial sample size}
  \DataTypeTok{matrix}\NormalTok{[N, P] X;}
  \DataTypeTok{vector}\NormalTok{[N] y;}
  \DataTypeTok{array}\NormalTok{[N] }\DataTypeTok{int}\NormalTok{\textless{}}\KeywordTok{lower}\NormalTok{=}\DecValTok{1}\NormalTok{, }\KeywordTok{upper}\NormalTok{=S\textgreater{} study;}
  \DataTypeTok{array}\NormalTok{[N] }\DataTypeTok{int}\NormalTok{\textless{}}\KeywordTok{lower}\NormalTok{=}\DecValTok{0}\NormalTok{, }\KeywordTok{upper}\NormalTok{=}\DecValTok{1}\NormalTok{\textgreater{} current;}
  \DataTypeTok{array}\NormalTok{[N] }\DataTypeTok{int}\NormalTok{\textless{}}\KeywordTok{lower}\NormalTok{=}\DecValTok{0}\NormalTok{, }\KeywordTok{upper}\NormalTok{=}\DecValTok{1}\NormalTok{\textgreater{} A;}
  \DataTypeTok{matrix}\NormalTok{[Ncur, P] Xcur;}
\NormalTok{\}}
\KeywordTok{parameters}\NormalTok{ \{}
  \DataTypeTok{real}\NormalTok{ mu\_alpha;}
  \DataTypeTok{real}\NormalTok{\textless{}}\KeywordTok{lower}\NormalTok{=}\DecValTok{0}\NormalTok{\textgreater{} sigma\_alpha;}
  \DataTypeTok{vector}\NormalTok{[S] z\_alpha;}
  \DataTypeTok{real}\NormalTok{ tau;}
  \DataTypeTok{vector}\NormalTok{[P] beta;}
  \DataTypeTok{real}\NormalTok{\textless{}}\KeywordTok{lower}\NormalTok{=}\DecValTok{0}\NormalTok{\textgreater{} sigma;}
\NormalTok{\}}
\KeywordTok{transformed parameters}\NormalTok{ \{}
  \DataTypeTok{vector}\NormalTok{[S] alpha;}
\NormalTok{  alpha = mu\_alpha + sigma\_alpha * z\_alpha;}
\NormalTok{\}}
\KeywordTok{model}\NormalTok{ \{}
\NormalTok{  mu\_alpha    \textasciitilde{} normal(}\DecValTok{0}\NormalTok{, }\DecValTok{2}\NormalTok{);}
\NormalTok{  sigma\_alpha \textasciitilde{} normal(}\DecValTok{0}\NormalTok{, }\DecValTok{1}\NormalTok{);}
\NormalTok{  z\_alpha     \textasciitilde{} normal(}\DecValTok{0}\NormalTok{, }\DecValTok{1}\NormalTok{);}
\NormalTok{  tau         \textasciitilde{} normal(}\DecValTok{0}\NormalTok{, }\DecValTok{2}\NormalTok{);}
\NormalTok{  beta        \textasciitilde{} normal(}\DecValTok{0}\NormalTok{, }\DecValTok{2}\NormalTok{);}
\NormalTok{  sigma       \textasciitilde{} normal(}\DecValTok{0}\NormalTok{, }\FloatTok{1.5}\NormalTok{);}

  \ControlFlowTok{for}\NormalTok{ (n }\ControlFlowTok{in} \DecValTok{1}\NormalTok{:N) \{}
\NormalTok{    y[n] \textasciitilde{} normal(alpha[study[n]] + tau * current[n] * A[n] + X[n] * beta, sigma);}
\NormalTok{  \}}
\NormalTok{\}}
\KeywordTok{generated quantities}\NormalTok{ \{}
  \DataTypeTok{vector}\NormalTok{[Ncur] mu1;}
  \DataTypeTok{vector}\NormalTok{[Ncur] mu0;}
  \DataTypeTok{real}\NormalTok{ muT;}
  \DataTypeTok{real}\NormalTok{ muC;}
  \DataTypeTok{real}\NormalTok{ Delta;}

\NormalTok{  mu1 = alpha[S] + tau * rep\_vector(}\DecValTok{1}\NormalTok{, Ncur) + Xcur * beta;}
\NormalTok{  mu0 = alpha[S] + Xcur * beta;}
\NormalTok{  muT = mean(mu1);}
\NormalTok{  muC = mean(mu0);}
\NormalTok{  Delta = muT {-} muC;}
\NormalTok{\}}
\end{Highlighting}
\end{Shaded}

\begin{center}\rule{0.5\linewidth}{0.5pt}\end{center}

\subsection{15. Stan Code: AR(1) Time-Dependent Borrowing
Model}\label{stan-code-ar1-time-dependent-borrowing-model}

\begin{Shaded}
\begin{Highlighting}[]
\KeywordTok{data}\NormalTok{ \{}
  \DataTypeTok{int}\NormalTok{\textless{}}\KeywordTok{lower}\NormalTok{=}\DecValTok{1}\NormalTok{\textgreater{} N;}
  \DataTypeTok{int}\NormalTok{\textless{}}\KeywordTok{lower}\NormalTok{=}\DecValTok{1}\NormalTok{\textgreater{} P;}
  \DataTypeTok{int}\NormalTok{\textless{}}\KeywordTok{lower}\NormalTok{=}\DecValTok{1}\NormalTok{\textgreater{} S;                 }\CommentTok{// ordered studies, S = 5}
  \DataTypeTok{int}\NormalTok{\textless{}}\KeywordTok{lower}\NormalTok{=}\DecValTok{1}\NormalTok{\textgreater{} Ncur;}
  \DataTypeTok{matrix}\NormalTok{[N, P] X;}
  \DataTypeTok{vector}\NormalTok{[N] y;}
  \DataTypeTok{array}\NormalTok{[N] }\DataTypeTok{int}\NormalTok{\textless{}}\KeywordTok{lower}\NormalTok{=}\DecValTok{1}\NormalTok{, }\KeywordTok{upper}\NormalTok{=S\textgreater{} study;}
  \DataTypeTok{array}\NormalTok{[N] }\DataTypeTok{int}\NormalTok{\textless{}}\KeywordTok{lower}\NormalTok{=}\DecValTok{0}\NormalTok{, }\KeywordTok{upper}\NormalTok{=}\DecValTok{1}\NormalTok{\textgreater{} current;}
  \DataTypeTok{array}\NormalTok{[N] }\DataTypeTok{int}\NormalTok{\textless{}}\KeywordTok{lower}\NormalTok{=}\DecValTok{0}\NormalTok{, }\KeywordTok{upper}\NormalTok{=}\DecValTok{1}\NormalTok{\textgreater{} A;}
  \DataTypeTok{matrix}\NormalTok{[Ncur, P] Xcur;}
\NormalTok{\}}
\KeywordTok{parameters}\NormalTok{ \{}
  \DataTypeTok{real}\NormalTok{ m0;}
  \DataTypeTok{real}\NormalTok{\textless{}}\KeywordTok{lower}\NormalTok{=}\DecValTok{0}\NormalTok{\textgreater{} s0;}
  \DataTypeTok{real}\NormalTok{ eta\_phi;}
  \DataTypeTok{real}\NormalTok{\textless{}}\KeywordTok{lower}\NormalTok{=}\DecValTok{0}\NormalTok{\textgreater{} omega;}
  \DataTypeTok{vector}\NormalTok{[S] alpha;}
  \DataTypeTok{real}\NormalTok{ tau;}
  \DataTypeTok{vector}\NormalTok{[P] beta;}
  \DataTypeTok{real}\NormalTok{\textless{}}\KeywordTok{lower}\NormalTok{=}\DecValTok{0}\NormalTok{\textgreater{} sigma;}
\NormalTok{\}}
\KeywordTok{transformed parameters}\NormalTok{ \{}
  \DataTypeTok{real}\NormalTok{ phi;}
\NormalTok{  phi = }\DecValTok{2}\NormalTok{ * inv\_logit(eta\_phi) {-} }\DecValTok{1}\NormalTok{;}
\NormalTok{\}}
\KeywordTok{model}\NormalTok{ \{}
\NormalTok{  m0    \textasciitilde{} normal(}\DecValTok{0}\NormalTok{, }\DecValTok{2}\NormalTok{);}
\NormalTok{  s0    \textasciitilde{} normal(}\DecValTok{0}\NormalTok{, }\DecValTok{1}\NormalTok{);}
\NormalTok{  eta\_phi \textasciitilde{} normal(}\DecValTok{0}\NormalTok{, }\DecValTok{1}\NormalTok{);}
\NormalTok{  omega \textasciitilde{} normal(}\DecValTok{0}\NormalTok{, }\DecValTok{1}\NormalTok{);}
\NormalTok{  tau   \textasciitilde{} normal(}\DecValTok{0}\NormalTok{, }\DecValTok{2}\NormalTok{);}
\NormalTok{  beta  \textasciitilde{} normal(}\DecValTok{0}\NormalTok{, }\DecValTok{2}\NormalTok{);}
\NormalTok{  sigma \textasciitilde{} normal(}\DecValTok{0}\NormalTok{, }\FloatTok{1.5}\NormalTok{);}

\NormalTok{  alpha[}\DecValTok{1}\NormalTok{] \textasciitilde{} normal(m0, s0);}
  \ControlFlowTok{for}\NormalTok{ (s }\ControlFlowTok{in} \DecValTok{2}\NormalTok{:S) \{}
\NormalTok{    alpha[s] \textasciitilde{} normal(phi * alpha[s {-} }\DecValTok{1}\NormalTok{], omega);}
\NormalTok{  \}}

  \ControlFlowTok{for}\NormalTok{ (n }\ControlFlowTok{in} \DecValTok{1}\NormalTok{:N) \{}
\NormalTok{    y[n] \textasciitilde{} normal(alpha[study[n]] + tau * current[n] * A[n] + X[n] * beta, sigma);}
\NormalTok{  \}}
\NormalTok{\}}
\KeywordTok{generated quantities}\NormalTok{ \{}
  \DataTypeTok{vector}\NormalTok{[Ncur] mu1;}
  \DataTypeTok{vector}\NormalTok{[Ncur] mu0;}
  \DataTypeTok{real}\NormalTok{ muT;}
  \DataTypeTok{real}\NormalTok{ muC;}
  \DataTypeTok{real}\NormalTok{ Delta;}

\NormalTok{  mu1 = alpha[S] + tau * rep\_vector(}\DecValTok{1}\NormalTok{, Ncur) + Xcur * beta;}
\NormalTok{  mu0 = alpha[S] + Xcur * beta;}
\NormalTok{  muT = mean(mu1);}
\NormalTok{  muC = mean(mu0);}
\NormalTok{  Delta = muT {-} muC;}
\NormalTok{\}}
\end{Highlighting}
\end{Shaded}

This is the main Stan model for time-dependent borrowing.

\begin{center}\rule{0.5\linewidth}{0.5pt}\end{center}

\subsection{16. Stan Code: Fixed Power Prior
Model}\label{stan-code-fixed-power-prior-model}

The simplest implementation uses a common intercept with weighted
historical likelihood contributions. Below, \texttt{a\_hist{[}h{]}}
contains the fixed power weights.

\begin{Shaded}
\begin{Highlighting}[]
\KeywordTok{data}\NormalTok{ \{}
  \DataTypeTok{int}\NormalTok{\textless{}}\KeywordTok{lower}\NormalTok{=}\DecValTok{1}\NormalTok{\textgreater{} N;}
  \DataTypeTok{int}\NormalTok{\textless{}}\KeywordTok{lower}\NormalTok{=}\DecValTok{1}\NormalTok{\textgreater{} P;}
  \DataTypeTok{int}\NormalTok{\textless{}}\KeywordTok{lower}\NormalTok{=}\DecValTok{1}\NormalTok{\textgreater{} H;}
  \DataTypeTok{matrix}\NormalTok{[N, P] X;}
  \DataTypeTok{vector}\NormalTok{[N] y;}
  \DataTypeTok{array}\NormalTok{[N] }\DataTypeTok{int}\NormalTok{\textless{}}\KeywordTok{lower}\NormalTok{=}\DecValTok{0}\NormalTok{, }\KeywordTok{upper}\NormalTok{=}\DecValTok{1}\NormalTok{\textgreater{} current;}
  \DataTypeTok{array}\NormalTok{[N] }\DataTypeTok{int}\NormalTok{\textless{}}\KeywordTok{lower}\NormalTok{=}\DecValTok{0}\NormalTok{, }\KeywordTok{upper}\NormalTok{=}\DecValTok{1}\NormalTok{\textgreater{} A;}
  \DataTypeTok{array}\NormalTok{[N] }\DataTypeTok{int}\NormalTok{\textless{}}\KeywordTok{lower}\NormalTok{=}\DecValTok{0}\NormalTok{, }\KeywordTok{upper}\NormalTok{=H\textgreater{} hist\_id; }\CommentTok{// 0 if current study, 1..H if historical}
  \DataTypeTok{vector}\NormalTok{\textless{}}\KeywordTok{lower}\NormalTok{=}\DecValTok{0}\NormalTok{, }\KeywordTok{upper}\NormalTok{=}\DecValTok{1}\NormalTok{\textgreater{}[H] a\_hist;}
\NormalTok{\}}
\KeywordTok{parameters}\NormalTok{ \{}
  \DataTypeTok{real}\NormalTok{ alpha;}
  \DataTypeTok{real}\NormalTok{ tau;}
  \DataTypeTok{vector}\NormalTok{[P] beta;}
  \DataTypeTok{real}\NormalTok{\textless{}}\KeywordTok{lower}\NormalTok{=}\DecValTok{0}\NormalTok{\textgreater{} sigma;}
\NormalTok{\}}
\KeywordTok{model}\NormalTok{ \{}
\NormalTok{  alpha \textasciitilde{} normal(}\DecValTok{0}\NormalTok{, }\DecValTok{2}\NormalTok{);}
\NormalTok{  tau   \textasciitilde{} normal(}\DecValTok{0}\NormalTok{, }\DecValTok{2}\NormalTok{);}
\NormalTok{  beta  \textasciitilde{} normal(}\DecValTok{0}\NormalTok{, }\DecValTok{2}\NormalTok{);}
\NormalTok{  sigma \textasciitilde{} normal(}\DecValTok{0}\NormalTok{, }\FloatTok{1.5}\NormalTok{);}

  \ControlFlowTok{for}\NormalTok{ (n }\ControlFlowTok{in} \DecValTok{1}\NormalTok{:N) \{}
    \ControlFlowTok{if}\NormalTok{ (current[n] == }\DecValTok{1}\NormalTok{) \{}
      \KeywordTok{target +=}\NormalTok{ normal\_lpdf(y[n] | alpha + tau * A[n] + X[n] * beta, sigma);}
\NormalTok{    \} }\ControlFlowTok{else}\NormalTok{ \{}
      \KeywordTok{target +=}\NormalTok{ a\_hist[hist\_id[n]] * normal\_lpdf(y[n] | alpha + X[n] * beta, sigma);}
\NormalTok{    \}}
\NormalTok{  \}}
\NormalTok{\}}
\end{Highlighting}
\end{Shaded}

A more advanced version could allow study-specific intercepts together
with weighted historical likelihoods.

\begin{center}\rule{0.5\linewidth}{0.5pt}\end{center}

\subsection{17. R Simulation Skeleton Using
cmdstanr}\label{r-simulation-skeleton-using-cmdstanr}

\begin{Shaded}
\begin{Highlighting}[]
\FunctionTok{library}\NormalTok{(cmdstanr)}
\FunctionTok{library}\NormalTok{(posterior)}
\FunctionTok{library}\NormalTok{(dplyr)}

\NormalTok{simulate\_one\_dataset }\OtherTok{\textless{}{-}} \ControlFlowTok{function}\NormalTok{(}
  \AttributeTok{n\_hist =} \FunctionTok{c}\NormalTok{(}\DecValTok{80}\NormalTok{, }\DecValTok{80}\NormalTok{, }\DecValTok{80}\NormalTok{, }\DecValTok{80}\NormalTok{),}
  \AttributeTok{n\_cur\_c =} \DecValTok{100}\NormalTok{,}
  \AttributeTok{n\_cur\_t =} \DecValTok{100}\NormalTok{,}
  \AttributeTok{beta =} \FunctionTok{c}\NormalTok{(}\FloatTok{0.6}\NormalTok{, }\SpecialCharTok{{-}}\FloatTok{0.5}\NormalTok{, }\FloatTok{0.4}\NormalTok{, }\SpecialCharTok{{-}}\FloatTok{0.3}\NormalTok{),}
  \AttributeTok{tau =} \SpecialCharTok{{-}}\FloatTok{0.3}\NormalTok{,}
  \AttributeTok{sigma =} \DecValTok{1}\NormalTok{,}
  \AttributeTok{alpha =} \FunctionTok{c}\NormalTok{(}\DecValTok{0}\NormalTok{, }\FloatTok{0.15}\NormalTok{, }\FloatTok{0.30}\NormalTok{, }\FloatTok{0.45}\NormalTok{, }\FloatTok{0.60}\NormalTok{)}
\NormalTok{) \{}
\NormalTok{  make\_x }\OtherTok{\textless{}{-}} \ControlFlowTok{function}\NormalTok{(n) \{}
    \FunctionTok{cbind}\NormalTok{(}
      \AttributeTok{x1 =} \FunctionTok{rnorm}\NormalTok{(n),}
      \AttributeTok{x2 =} \FunctionTok{rbinom}\NormalTok{(n, }\DecValTok{1}\NormalTok{, }\FloatTok{0.5}\NormalTok{),}
      \AttributeTok{x3 =} \FunctionTok{rnorm}\NormalTok{(n),}
      \AttributeTok{x4 =} \FunctionTok{rbinom}\NormalTok{(n, }\DecValTok{1}\NormalTok{, }\FloatTok{0.4}\NormalTok{)}
\NormalTok{    )}
\NormalTok{  \}}

\NormalTok{  dat\_list }\OtherTok{\textless{}{-}} \FunctionTok{list}\NormalTok{()}

  \ControlFlowTok{for}\NormalTok{ (h }\ControlFlowTok{in} \DecValTok{1}\SpecialCharTok{:}\DecValTok{4}\NormalTok{) \{}
\NormalTok{    Xh }\OtherTok{\textless{}{-}} \FunctionTok{make\_x}\NormalTok{(n\_hist[h])}
\NormalTok{    yh }\OtherTok{\textless{}{-}}\NormalTok{ alpha[h] }\SpecialCharTok{+}\NormalTok{ Xh }\SpecialCharTok{\%*\%}\NormalTok{ beta }\SpecialCharTok{+} \FunctionTok{rnorm}\NormalTok{(n\_hist[h], }\DecValTok{0}\NormalTok{, sigma)}
\NormalTok{    dat\_list[[h]] }\OtherTok{\textless{}{-}} \FunctionTok{data.frame}\NormalTok{(}
      \AttributeTok{study =}\NormalTok{ h,}
      \AttributeTok{current =} \DecValTok{0}\NormalTok{,}
      \AttributeTok{A =} \DecValTok{0}\NormalTok{,}
      \AttributeTok{y =} \FunctionTok{as.numeric}\NormalTok{(yh),}
\NormalTok{      Xh}
\NormalTok{    )}
\NormalTok{  \}}

\NormalTok{  Xc0 }\OtherTok{\textless{}{-}} \FunctionTok{make\_x}\NormalTok{(n\_cur\_c)}
\NormalTok{  yc0 }\OtherTok{\textless{}{-}}\NormalTok{ alpha[}\DecValTok{5}\NormalTok{] }\SpecialCharTok{+}\NormalTok{ Xc0 }\SpecialCharTok{\%*\%}\NormalTok{ beta }\SpecialCharTok{+} \FunctionTok{rnorm}\NormalTok{(n\_cur\_c, }\DecValTok{0}\NormalTok{, sigma)}
\NormalTok{  d0 }\OtherTok{\textless{}{-}} \FunctionTok{data.frame}\NormalTok{(}\AttributeTok{study =} \DecValTok{5}\NormalTok{, }\AttributeTok{current =} \DecValTok{1}\NormalTok{, }\AttributeTok{A =} \DecValTok{0}\NormalTok{, }\AttributeTok{y =} \FunctionTok{as.numeric}\NormalTok{(yc0), Xc0)}

\NormalTok{  Xc1 }\OtherTok{\textless{}{-}} \FunctionTok{make\_x}\NormalTok{(n\_cur\_t)}
\NormalTok{  yc1 }\OtherTok{\textless{}{-}}\NormalTok{ alpha[}\DecValTok{5}\NormalTok{] }\SpecialCharTok{+}\NormalTok{ tau }\SpecialCharTok{+}\NormalTok{ Xc1 }\SpecialCharTok{\%*\%}\NormalTok{ beta }\SpecialCharTok{+} \FunctionTok{rnorm}\NormalTok{(n\_cur\_t, }\DecValTok{0}\NormalTok{, sigma)}
\NormalTok{  d1 }\OtherTok{\textless{}{-}} \FunctionTok{data.frame}\NormalTok{(}\AttributeTok{study =} \DecValTok{5}\NormalTok{, }\AttributeTok{current =} \DecValTok{1}\NormalTok{, }\AttributeTok{A =} \DecValTok{1}\NormalTok{, }\AttributeTok{y =} \FunctionTok{as.numeric}\NormalTok{(yc1), Xc1)}

\NormalTok{  dat }\OtherTok{\textless{}{-}} \FunctionTok{bind\_rows}\NormalTok{(dat\_list, }\FunctionTok{list}\NormalTok{(d0, d1))}
\NormalTok{  dat}
\NormalTok{\}}

\NormalTok{prepare\_stan\_data\_ar1 }\OtherTok{\textless{}{-}} \ControlFlowTok{function}\NormalTok{(dat) \{}
\NormalTok{  X }\OtherTok{\textless{}{-}} \FunctionTok{as.matrix}\NormalTok{(dat[, }\FunctionTok{c}\NormalTok{(}\StringTok{"x1"}\NormalTok{, }\StringTok{"x2"}\NormalTok{, }\StringTok{"x3"}\NormalTok{, }\StringTok{"x4"}\NormalTok{)])}
\NormalTok{  cur }\OtherTok{\textless{}{-}} \FunctionTok{subset}\NormalTok{(dat, study }\SpecialCharTok{==} \DecValTok{5}\NormalTok{)}
\NormalTok{  Xcur }\OtherTok{\textless{}{-}} \FunctionTok{as.matrix}\NormalTok{(cur[, }\FunctionTok{c}\NormalTok{(}\StringTok{"x1"}\NormalTok{, }\StringTok{"x2"}\NormalTok{, }\StringTok{"x3"}\NormalTok{, }\StringTok{"x4"}\NormalTok{)])}

  \FunctionTok{list}\NormalTok{(}
    \AttributeTok{N =} \FunctionTok{nrow}\NormalTok{(dat),}
    \AttributeTok{P =} \FunctionTok{ncol}\NormalTok{(X),}
    \AttributeTok{S =} \FunctionTok{length}\NormalTok{(}\FunctionTok{unique}\NormalTok{(dat}\SpecialCharTok{$}\NormalTok{study)),}
    \AttributeTok{Ncur =} \FunctionTok{nrow}\NormalTok{(Xcur),}
    \AttributeTok{X =}\NormalTok{ X,}
    \AttributeTok{y =}\NormalTok{ dat}\SpecialCharTok{$}\NormalTok{y,}
    \AttributeTok{study =}\NormalTok{ dat}\SpecialCharTok{$}\NormalTok{study,}
    \AttributeTok{current =}\NormalTok{ dat}\SpecialCharTok{$}\NormalTok{current,}
    \AttributeTok{A =}\NormalTok{ dat}\SpecialCharTok{$}\NormalTok{A,}
    \AttributeTok{Xcur =}\NormalTok{ Xcur}
\NormalTok{  )}
\NormalTok{\}}

\NormalTok{extract\_estimands }\OtherTok{\textless{}{-}} \ControlFlowTok{function}\NormalTok{(fit) \{}
\NormalTok{  draws }\OtherTok{\textless{}{-}}\NormalTok{ fit}\SpecialCharTok{$}\FunctionTok{draws}\NormalTok{(}\FunctionTok{c}\NormalTok{(}\StringTok{"muT"}\NormalTok{, }\StringTok{"muC"}\NormalTok{, }\StringTok{"Delta"}\NormalTok{))}
  \FunctionTok{as\_draws\_df}\NormalTok{(draws)}
\NormalTok{\}}
\end{Highlighting}
\end{Shaded}

\begin{center}\rule{0.5\linewidth}{0.5pt}\end{center}

\subsection{18. Planned Output Tables and
Figures}\label{planned-output-tables-and-figures}

The final simulation report should contain:

\begin{enumerate}
\def\labelenumi{\arabic{enumi}.}
\tightlist
\item
  A table summarizing all simulation scenarios.
\item
  A table of true parameter values used for each scenario.
\item
  Bias / RMSE / coverage tables for (\mu\_T), (\mu\_C), and (\Delta).
\item
  A plot of empirical bias against drift severity.
\item
  A plot of interval width against drift severity.
\item
  Power and type I error plots by method.
\item
  Posterior mean trajectories of study intercepts under the AR(1) model.
\end{enumerate}

\begin{center}\rule{0.5\linewidth}{0.5pt}\end{center}

\subsection{19. Interpretation Strategy}\label{interpretation-strategy}

The interpretation should emphasize:

\begin{itemize}
\tightlist
\item
  whether explicit time modeling protects against drift-induced bias,
\item
  when full borrowing is efficient but unsafe,
\item
  whether exchangeable hierarchical priors are sufficiently robust when
  drift exists,
\item
  how much efficiency is gained relative to no borrowing,
\item
  and whether posterior g-computation recovers the marginal estimand
  under covariate imbalance.
\end{itemize}

\begin{center}\rule{0.5\linewidth}{0.5pt}\end{center}

\subsection{20. Recommended Primary Analysis for the
Thesis}\label{recommended-primary-analysis-for-the-thesis}

For the main thesis simulation, a strong primary design would be:

\begin{itemize}
\tightlist
\item
  4 historical studies plus 1 current trial,
\item
  covariate-adjusted Gaussian outcome model,
\item
  estimands (\mu\_T), (\mu\_C), and (\Delta) for the current population,
\item
  methods compared: no borrowing, full borrowing, hierarchical prior,
  AR(1) prior, and power prior,
\item
  primary drift scenarios: none, linear drift, and AR(1) drift,
\item
  main treatment settings: (\tau=0) and (\tau=-0.3).
\end{itemize}

This provides a clean framework that is aligned with the methodological
contribution of time-dependent borrowing.

\begin{center}\rule{0.5\linewidth}{0.5pt}\end{center}

\subsection{21. Next Extension}\label{next-extension}

A natural next extension is to allow the regression coefficients for the
measured covariates to be common across studies while permitting the
intercept to vary over time. This separates:

\begin{itemize}
\tightlist
\item
  \textbf{measured prognostic adjustment} through (X\^{}\top\beta), and
\item
  \textbf{unmeasured temporal change} through the evolving intercept
  (\alpha\_j).
\end{itemize}

This is exactly the structure needed to represent evolving standard of
care or other unmeasured secular changes.

\begin{center}\rule{0.5\linewidth}{0.5pt}\end{center}

\subsection{22. Short Thesis-Style Summary
Paragraph}\label{short-thesis-style-summary-paragraph}

The simulation study will evaluate Bayesian borrowing methods for
continuous outcomes in randomized trials with multiple historical
control datasets. Patient-level covariates will be incorporated through
an outcome regression model, and the target estimands will be the
marginal treatment and control means, together with the marginal
treatment effect, in the current trial population. Historical and
current control information will be linked through alternative prior
structures, including exchangeable hierarchical priors, power priors,
and a time-dependent AR(1) prior on study-specific control intercepts.
Simulation scenarios will vary the degree of temporal drift, treatment
effect size, sample size, and covariate shift, and performance will be
assessed using bias, RMSE, interval coverage, interval width, power, and
type I error. The primary objective is to determine whether explicit
modeling of temporal dependence improves robustness and efficiency when
borrowing from historical controls.




\end{document}
